% Prose rendering of PrimeTensor/Analysis/Limit.lean.
% Fragment: no \documentclass, no preamble, no document environment.

\section{Intrinsic convergence on the completed carrier}
\label{Analysis:Limit:sec}

\leanfile{PrimeTensor/Analysis/Limit.lean}

\subsection*{What this file establishes}

$\name{MulReal}$ has so far been an algebraic object: a commutative carrier in
the sense of \texttt{PrimeTensor/Structure.lean}, receiving a
structure-preserving map from $\name{MulRat}$.  Nothing has related its
elements to the scale they were built from.  This file supplies the missing
analytic vocabulary --- a nearness relation on completed points and a notion of
convergence for sequences of them --- and proves the one theorem that makes the
construction deserve the name completion:
\[
  \text{for every Cauchy stream } a, \qquad
  \name{ofRat}(a_n) \longrightarrow [a] .
\]
That is, the embedded rational terms of a Cauchy stream converge to the point
that stream represents.  Read the other way, it says that
$\name{ofRat}(\name{MulRat})$ is dense in $\name{MulReal}$: every completed
point is a limit of embedded rationals, with an explicit sequence exhibiting it.

The design question the file has to answer is how to transport a relation
defined on representatives across a quotient whose relation is not itself the
one being transported.  Its answer --- quantify existentially over
representatives --- is what
Design note~\ref{Analysis:Limit:note:existential} is about, and it is what
makes the central theorem cost nothing.

\subsection{Nearness of completed points}

\begin{definition}[\lean{MulReal.ScaleNear}]\label{Analysis:Limit:def:scalenear}
Completed points $x$ and $y$ are \emph{near at level $\ell$} when there exist
Cauchy streams $a$ and $b$ with $[a] = x$ and $[b] = y$, and an anchor beyond
which their terms are $\name{ScaleWithin}(\ell,-,-)$.
\end{definition}

\begin{leancode}
def ScaleNear (level : Depth) (x y : MulReal) : Prop :=
  ∃ a b : MulCauchyStream,
    ofStream a = x ∧
    ofStream b = y ∧
    ∃ anchor : Depth,
      ∀ n : Depth,
        Depth.AtOrAfter anchor n →
        MulRat.ScaleWithin level (a.term n) (b.term n)
\end{leancode}

\begin{designnote}[Existential, not universal, over representatives]
\label{Analysis:Limit:note:existential}
A relation on representatives descends to a quotient when it is invariant
under the quotient's own relation.  $\name{ScaleWithin}$ at a fixed level is
not: two streams asymptotic in the sense of
\texttt{PrimeTensor/Completion.lean} agree only up to the composition law of
\texttt{PrimeTensor/Scale/Laws.lean}, and that law consumes a level each time
it is used.  Replacing representatives therefore cannot be free at a fixed
level, and a relation defined by ``\emph{for all} representatives'' and one
defined by ``\emph{there exist} representatives'' are genuinely different
propositions here, the second being the weaker.

The file takes the weaker one deliberately --- the docstring calls the choice
part of the quotient semantics.  The consequence, visible immediately below, is
that establishing $\name{ScaleNear}$ means exhibiting one convenient pair of
representatives rather than controlling every pair, and the central theorem
gets to choose the pair that makes it trivial.  The price is paid in what can
be proved \emph{about} the relation, and
Remark~\ref{Analysis:Limit:rem:not-transitive} is the bill.

Note also what is not introduced.  There is no metric, no distance function,
and no additive $\varepsilon$: nearness is a family of relations indexed by
depth, exactly as at the $\name{MulRat}$ layer, and the index is the same
$\name{Depth}$ that indexes the sequences.
\end{designnote}

\begin{lemma}[\lean{MulReal.scaleNear\_refl} and \lean{scaleNear\_symm}]
\label{Analysis:Limit:lem:refl-symm}
$\name{ScaleNear}(\ell, x, x)$ for every level and every $x$; and
$\name{ScaleNear}(\ell,x,y)$ implies $\name{ScaleNear}(\ell,y,x)$.
\end{lemma}

\begin{proof}
For reflexivity, induct to choose a representative $a$ of $x$ and offer the
pair $(a,a)$, with anchor $\ctor{one}$ so the tail is everything; each term
comparison is $\name{scaleWithin\_refl}$.  For symmetry, take the witnessing
pair supplied by the hypothesis, exchange its two components together with the
two identifications, keep the anchor and the level, and apply
$\name{scaleWithin\_symm}$ termwise.
\end{proof}

\begin{remark}[Transitivity is not proved, at any level]
\label{Analysis:Limit:rem:not-transitive}
$\name{ScaleWithin}$ was not transitive at a fixed level either, but it had the
level-shifted substitute of \texttt{PrimeTensor/Scale/Laws.lean}.  No analogue
is proved here, and the obstruction is exactly the existential of
Design note~\ref{Analysis:Limit:note:existential}: a witness for
$\name{ScaleNear}(\ell,x,y)$ names some representative of $y$, a witness for
$\name{ScaleNear}(\ell,y,z)$ names \emph{some other} one, and the two need not
be the same stream.  Chaining them requires first crossing from one
representative of $y$ to the other, which is possible --- they are asymptotic,
so they are eventually close at every level --- but costs a further application
of the composition law, and then a comparison between two different levels,
which is the monotonicity statement that
\texttt{PrimeTensor/Scale/Laws.lean} left unproved.  So a level-shifted
transitivity is plausible and is not available off the shelf.  Nothing below
needs it.
\end{remark}

\subsection{Convergence}

\begin{definition}[\lean{MulReal.Seq} and \lean{MulReal.ConvergesTo}]
\label{Analysis:Limit:def:converges}
A \emph{sequence} is a function $s : \name{Depth} \to \name{MulReal}$, and
$s$ \emph{converges to} $x$ when
\[
  \forall\, \ell\; \exists\, c\; \forall\, n :
    \name{AtOrAfter}(c,n) \longrightarrow \name{ScaleNear}(\ell, s_n, x).
\]
\end{definition}

\begin{leancode}
def ConvergesTo (s : Seq) (x : MulReal) : Prop :=
  ∀ level : Depth,
    ∃ anchor : Depth,
      ∀ n : Depth,
        Depth.AtOrAfter anchor n →
        ScaleNear level (s n) x
\end{leancode}

The shape is the one used twice already --- for $\name{IsMulCauchy}$ and for
$\name{MulAsymptotic}$ in \texttt{PrimeTensor/Completion.lean} --- with
$\name{ScaleNear}$ in place of $\name{ScaleWithin}$: for every level a tail on
which the requirement holds, with the level and the tail anchor both depths and
nothing else quantified.

\begin{lemma}[\lean{MulReal.converges\_constant}]\label{Analysis:Limit:lem:constant}
The constant sequence at $x$ converges to $x$.
\end{lemma}

\begin{proof}
Anchor $\ctor{one}$, and each requirement is
Lemma~\ref{Analysis:Limit:lem:refl-symm}.
\end{proof}

\begin{lemma}[\lean{MulReal.ofStream\_eq\_of\_asymptotic}]
\label{Analysis:Limit:lem:sound}
Asymptotic streams represent the same completed point.
\end{lemma}

\begin{proof}
This is $\ctor{Quotient.sound}$ for the setoid of
\texttt{PrimeTensor/Completion.lean}, named so that later proofs can rewrite
along it.
\end{proof}

\subsection{The convergence theorem}

\begin{theorem}[\lean{MulReal.terms\_converge}]\label{Analysis:Limit:thm:main}
For every Cauchy stream $a$, the sequence $n \mapsto \name{ofRat}(a_n)$
converges to $[a]$.
\end{theorem}

\begin{leancode}
theorem terms_converge (a : MulCauchyStream) :
    ConvergesTo (fun n => ofRat (a.term n)) (ofStream a) := by
  intro level
  obtain ⟨anchor, hcauchy⟩ := a.cauchy level
  refine ⟨anchor, ?_⟩
  intro n hn
  refine ⟨MulCauchyStream.constant (a.term n), a, rfl, rfl, anchor, ?_⟩
  intro k hk
  exact hcauchy n k hn hk
\end{leancode}

\begin{proof}
Fix a level $\ell$ and let $c$ be the anchor supplied by the Cauchy property of
$a$ \emph{at that same level} --- not at $\ctor{succ}\,\ell$.  Take $c$ as the
anchor for the convergence claim as well.

Let $n$ be at or after $c$.  To witness
$\name{ScaleNear}(\ell, \name{ofRat}(a_n), [a])$, offer the pair consisting of
the constant stream at $a_n$ and the stream $a$ itself.  The two
identifications are $\ctor{rfl}$: $\name{ofRat}(q)$ was defined in
\texttt{PrimeTensor/Completion.lean} as the class of the constant stream at
$q$, and the second component is $a$ unchanged.  Take $c$ once more as the
anchor of the witness.  For $k$ at or after $c$, the required comparison is
\[
  \name{ScaleWithin}\bigl(\ell,\; (\name{constant}(a_n))_k,\; a_k\bigr)
  \;=\;
  \name{ScaleWithin}(\ell,\, a_n,\, a_k),
\]
and both $n$ and $k$ lie in the tail at $c$, so this is precisely the Cauchy
property of $a$ at $\ell$ applied to the pair $(n,k)$.
\end{proof}

\begin{remark}[The proof costs nothing, and why]\label{Analysis:Limit:rem:free}
No level is spent: the Cauchy property is invoked at $\ell$ and the conclusion
is at $\ell$.  No $\name{join}$ of anchors appears; one anchor serves for the
convergence claim and for every witness.  The whole theorem is the Cauchy
property of $a$, re-read.

That is the payoff of
Design note~\ref{Analysis:Limit:note:existential}.  The witnessing pair may be
chosen, and the choice made here --- compare the constant stream at $a_n$
against $a$ itself, rather than against some other representative --- turns the
statement ``the $n$-th embedded term is near the limit'' into the statement
``$a_n$ is close to $a_k$ for large $k$'', which is what being Cauchy says.
Under a universally quantified $\name{ScaleNear}$ the same argument would have
to survive an arbitrary change of representative on both sides, and would spend
levels doing it.
\end{remark}

\begin{corollary}[\lean{MulReal.terms\_converge\_of\_asymptotic}]
\label{Analysis:Limit:cor:asymptotic}
If $a$ and $b$ are asymptotic then $n \mapsto \name{ofRat}(a_n)$ converges to
$[b]$.
\end{corollary}

\begin{proof}
$[a] = [b]$ by Lemma~\ref{Analysis:Limit:lem:sound}; rewrite the target and
apply Theorem~\ref{Analysis:Limit:thm:main}.
\end{proof}

\begin{remark}[Density, and the half of completeness this is]
\label{Analysis:Limit:rem:density}
Theorem~\ref{Analysis:Limit:thm:main} says every completed point is a limit of
embedded rationals, and it does more than assert existence: the approximating
sequence is exhibited, being the terms of any representative.  That is the
density half of the statement that $\name{MulReal}$ completes $\name{MulRat}$.

The other half --- that $\name{MulReal}$ is itself complete, every Cauchy
sequence of completed points converging in $\name{MulReal}$ --- is not proved,
and cannot yet be stated: $\name{IsMulCauchy}$ is a predicate on
$\name{MulRat}$-valued streams, and no Cauchy condition for a
$\name{MulReal}$-valued $\name{Seq}$ is defined anywhere in the file.  The
vocabulary for the second half is missing, not merely its proof.
\end{remark}

\begin{remark}[Limits are not yet unique]\label{Analysis:Limit:rem:unique}
$\name{ConvergesTo}$ is a relation, and nothing here makes it a function.  If
$s$ converges to $x$ and to $y$, no lemma concludes $x = y$; the usual argument
would combine the two through a transitivity that
Remark~\ref{Analysis:Limit:rem:not-transitive} says is unavailable, and would
then need the separation statement that two points near at every level are
equal.  That statement is the same missing ingredient as the injectivity of
$\name{ofRat}$, noted at the end of \texttt{PrimeTensor/Structure.lean}: both
are the assertion that the ladder of levels actually separates points.  Until
it is proved one may not speak of \emph{the} limit of a sequence, and this file
never does.
\end{remark}

\subsection*{Summary of the correspondence}

\begin{center}
\small
\begin{tabular}{@{}lll@{}}
\hline
Lean declaration & Kind & Here \\
\hline
\lean{MulReal.ScaleNear}          & definition & Def.~\ref{Analysis:Limit:def:scalenear} \\
\lean{MulReal.scaleNear\_refl}    & theorem    & Lem.~\ref{Analysis:Limit:lem:refl-symm} \\
\lean{MulReal.scaleNear\_symm}    & theorem    & Lem.~\ref{Analysis:Limit:lem:refl-symm} \\
\lean{MulReal.Seq}                & abbrev     & Def.~\ref{Analysis:Limit:def:converges} \\
\lean{MulReal.ConvergesTo}        & definition & Def.~\ref{Analysis:Limit:def:converges} \\
\lean{MulReal.converges\_constant} & theorem   & Lem.~\ref{Analysis:Limit:lem:constant} \\
\lean{MulReal.ofStream\_eq\_of\_asymptotic} & theorem & Lem.~\ref{Analysis:Limit:lem:sound} \\
\lean{MulReal.terms\_converge}    & theorem    & Thm.~\ref{Analysis:Limit:thm:main} \\
\lean{MulReal.terms\_converge\_of\_asymptotic} & theorem & Cor.~\ref{Analysis:Limit:cor:asymptotic} \\
\hline
\end{tabular}
\end{center}

\noindent
Every declaration of \texttt{PrimeTensor/Analysis/Limit.lean} appears above.
The file contains no \lean{sorry}, no axiom, and no unproved named proposition.
Remarks~\ref{Analysis:Limit:rem:not-transitive}, \ref{Analysis:Limit:rem:free},
\ref{Analysis:Limit:rem:density} and~\ref{Analysis:Limit:rem:unique} are stated
for the reader and are not declarations of the file.
